\chapter{Về Kỷ Luật}
\markboth{Về Kỷ Luật}{On Discipline}

\section{Kỷ luật là tự thương lấy tương lai mình.}
\begin{truyen}{Tự thương tương lai}{Discipline Protects Your Future}
\chuhoa{T}{hầy nói: ``Nhiều em tưởng kỷ luật là bị ép. Không. Kỷ luật là cách em tự thương lấy tương lai mình. Em ngủ đúng giờ, học đúng lúc, giữ lời với chính mình~--- tức là em đang đỡ cho chính mình khỏi phải trả giá sau này.''}
\textit{(The teacher said: ``Many students think discipline is pressure. It isn't. Discipline is how you protect your future self. When you sleep on time, study on time, and keep promises to yourself, you save yourself from paying a bigger price later.'')}
\end{truyen}

\begin{baihoc}
Kỷ luật không lấy mất tự do; nó giữ cho tương lai của em bớt hỗn loạn.
\textit{(Discipline doesn't steal freedom; it keeps your future from becoming chaotic.)}
\end{baihoc}

\begin{ghinhoanh}
\textbf{Short English to remember:}

Discipline is self-respect for your future.

\end{ghinhoanh}

\begin{tuhocnhanh}
1. Em đang giữ kỷ luật nào để giúp chính mình? \textit{What discipline are you keeping to help yourself?}

2. Nếu thiếu kỷ luật, tương lai em sẽ trả giá ở đâu? \textit{Where will your future pay the price if discipline is missing?}
\end{tuhocnhanh}
\ngancach

\section{Giờ giấc là bài học đầu tiên của người trưởng thành.}
\begin{truyen}{Đến đúng giờ}{The Lesson of Time}
\chuhoa{C}{ó em học không tệ, nhưng hay đến muộn. Thầy không mắng dài. Thầy chỉ nói: ``Giờ giấc là bài học đầu tiên của người trưởng thành. Người đến muộn nhiều lần đang nói với thế giới rằng thời gian của người khác không quan trọng.''}
\textit{(One student was capable but often late. The teacher didn't lecture long. He only said: ``Time is the first lesson of adulthood. Someone who is repeatedly late tells the world that other people's time doesn't matter.'')}
\end{truyen}

\begin{baihoc}
Đi đúng giờ là một biểu hiện nhỏ của sự tôn trọng và đáng tin.
\textit{(Being on time is a small sign of respect and reliability.)}
\end{baihoc}

\begin{ghinhoanh}
\textbf{Short English to remember:}

Punctuality is the first language of maturity.

\end{ghinhoanh}

\begin{tuhocnhanh}
1. Vì sao đi đúng giờ liên quan đến nhân cách? \textit{Why is punctuality related to character?}

2. Em cần sửa thói quen nào để bớt trễ? \textit{Which habit do you need to change to be less late?}
\end{tuhocnhanh}
\ngancach

\section{Làm việc khó trước khi làm việc thích.}
\begin{truyen}{Việc khó trước}{Hard Things First}
\chuhoa{T}{hầy bảo: ``Nếu em chỉ làm việc mình thích, em sẽ tiến rất chậm. Người biết đi xa là người làm việc khó trước, việc thích sau. Vì chính việc khó mới mở đường cho em lớn lên.''}
\textit{(The teacher said: ``If you only do what you like, you'll grow slowly. The people who go far do the hard things first and the pleasant things later. The hard things are what make you grow.'')}
\end{truyen}

\begin{baihoc}
Trưởng thành bắt đầu từ lúc em biết ưu tiên điều cần hơn điều thích.
\textit{(Growth begins when you choose what is needed over what is merely pleasant.)}
\end{baihoc}

\begin{ghinhoanh}
\textbf{Short English to remember:}

Do the hard thing first.

\end{ghinhoanh}

\begin{tuhocnhanh}
1. Việc khó nào em đang né tránh? \textit{Which hard task are you avoiding?}

2. Nếu làm nó trước, em sẽ nhẹ đi điều gì? \textit{What becomes lighter if you do it first?}
\end{tuhocnhanh}
\ngancach

\section{Tự nhắc mình còn quý hơn để người khác nhắc.}
\begin{truyen}{Tự nhắc}{Remind Yourself First}
\chuhoa{M}{ỗi lần lớp ồn, thầy nhắc một lần. Rồi thầy dừng lại và nói: ``Điều quý nhất không phải là có người nhắc em. Điều quý nhất là em tự biết lúc nào cần chỉnh mình trước khi người khác phải lên tiếng.''}
\textit{(Whenever the class got noisy, the teacher reminded them once. Then he stopped and said: ``The valuable thing is not that someone reminds you. The valuable thing is knowing when to correct yourself before someone has to say it.'')}
\end{truyen}

\begin{baihoc}
Tự giác là cấp độ cao hơn của vâng lời.
\textit{(Self-awareness is a higher level than simple obedience.)}
\end{baihoc}

\begin{ghinhoanh}
\textbf{Short English to remember:}

The best reminder comes from inside.

\end{ghinhoanh}

\begin{tuhocnhanh}
1. Em có thường chờ người khác nhắc mới làm không? \textit{Do you often wait for others to remind you?}

2. Em muốn rèn sự tự giác ở việc nào trước? \textit{Which area do you want to build self-discipline in first?}
\end{tuhocnhanh}
\ngancach

\section{Người đáng tin là người làm điều đã hứa.}
\begin{truyen}{Giữ lời}{Keeping Your Word}
\chuhoa{C}{ó em hứa nộp bài rồi lại quên. Thầy nói nhẹ: ``Người đáng tin không phải người nói hay. Người đáng tin là người làm điều đã hứa. Từ việc nhỏ như nộp bài đúng hẹn, người ta mới dần tin em ở việc lớn hơn.''}
\textit{(A student promised to submit work and then forgot. The teacher said gently: ``A trustworthy person is not the one who speaks well, but the one who does what was promised. Trust in big things starts from small things like handing in work on time.'')}
\end{truyen}

\begin{baihoc}
Giữ lời ở việc nhỏ là nền của uy tín ở việc lớn.
\textit{(Keeping small promises is the foundation of credibility in larger matters.)}
\end{baihoc}

\begin{ghinhoanh}
\textbf{Short English to remember:}

Trust grows where promises are kept.

\end{ghinhoanh}

\begin{tuhocnhanh}
1. Em đã hứa với ai điều gì mà còn chưa làm? \textit{What promise have you not yet fulfilled?}

2. Vì sao việc giữ lời quan trọng hơn lời nói đẹp? \textit{Why is keeping your word more important than speaking well?}
\end{tuhocnhanh}
\ngancach

\section{Mỗi ngày đúng giờ là một lần thắng bản thân.}
\begin{truyen}{Thắng mình trước}{Winning Over Yourself}
\chuhoa{T}{hầy nói: ``Nghe có vẻ nhỏ, nhưng mỗi ngày đúng giờ là một lần thắng bản thân. Em thắng cái lười thêm năm phút, thắng cái trì hoãn, thắng cái dễ dãi với chính mình. Người thắng mình đều đặn sẽ không dễ thua đời.''}
\textit{(The teacher said: ``It sounds small, but being on time each day is a victory over yourself. You beat the urge to sleep five more minutes, beat delay, beat being too easy on yourself. Someone who keeps winning over themselves won't easily lose to life.'')}
\end{truyen}

\begin{baihoc}
Những chiến thắng nhỏ trước bản thân tạo nên sức mạnh bền nhất.
\textit{(Small victories over yourself build the most durable strength.)}
\end{baihoc}

\begin{ghinhoanh}
\textbf{Short English to remember:}

Small self-victories build great strength.

\end{ghinhoanh}

\begin{tuhocnhanh}
1. Hôm nay em đã thắng mình ở điều gì? \textit{How did you win over yourself today?}

2. Nếu thắng bản thân thêm một việc mỗi ngày, em sẽ khác ra sao? \textit{How would you change if you beat one more weakness each day?}
\end{tuhocnhanh}
\ngancach

\section{Kỷ luật nhỏ giữ được ước mơ lớn.}
\begin{truyen}{Cái lớn được giữ bởi cái nhỏ}{Small Discipline, Big Dream}
\chuhoa{M}{ột em nói rằng em có ước mơ lớn. Thầy gật đầu rồi hỏi: ``Em có giữ được giờ ngủ, giờ học, giờ ôn bài không?'' Em ấy im lặng. Thầy bảo: ``Ước mơ lớn không chết vì thiếu lời nói đẹp; nó chết vì thiếu kỷ luật nhỏ mỗi ngày.''}
\textit{(A student said he had a big dream. The teacher nodded and asked: ``Can you keep your sleep time, study time, and review time?'' The student fell silent. The teacher said: ``Big dreams do not die from lack of beautiful words. They die from a lack of small daily discipline.'')}
\end{truyen}

\begin{baihoc}
Ước mơ bền không đứng trên cảm hứng, mà đứng trên nề nếp.
\textit{(A lasting dream stands not on emotion, but on routine.)}
\end{baihoc}

\begin{ghinhoanh}
\textbf{Short English to remember:}

Big dreams need small discipline.

\end{ghinhoanh}

\begin{tuhocnhanh}
1. Ước mơ nào của em đang cần một kỷ luật nhỏ đi kèm? \textit{Which dream of yours needs a small discipline to support it?}

2. Em có đang chỉ mơ mà chưa rèn nề nếp không? \textit{Are you dreaming without building routine?}
\end{tuhocnhanh}
\ngancach

\section{Khi không muốn làm mà vẫn làm, em đang lớn lên.}
\begin{truyen}{Lớn lên trong việc khó}{Growing Through Effort}
\chuhoa{T}{hầy nói: ``Lúc em muốn làm mà làm được, đó là bình thường. Lúc em không muốn làm mà vẫn làm~--- đó mới là lúc em lớn lên. Vì khi ấy em không còn bị cảm xúc kéo đi nữa.''}
\textit{(The teacher said: ``When you want to do something and do it, that's normal. When you do it even when you don't feel like it, that is when you grow. In that moment, your emotions no longer control you.'')}
\end{truyen}

\begin{baihoc}
Sự trưởng thành thật thường bắt đầu trong những việc em không hứng thú nhưng vẫn hoàn thành.
\textit{(Real maturity often begins in tasks you do not feel excited about but still complete.)}
\end{baihoc}

\begin{ghinhoanh}
\textbf{Short English to remember:}

You grow when you do it anyway.

\end{ghinhoanh}

\begin{tuhocnhanh}
1. Việc nào em không thích nhưng biết mình cần làm? \textit{Which task do you dislike but know you need to do?}

2. Em sẽ thử hoàn thành nó bằng cách nào hôm nay? \textit{How will you try to finish it today?}
\end{tuhocnhanh}
\ngancach

\section{Bàn học gọn cũng là một dạng tư duy rõ.}
\begin{truyen}{Gọn bên ngoài, sáng bên trong}{Clear Space, Clear Mind}
\chuhoa{T}{hầy đi qua một bàn học đầy giấy vụn, sách lộn xộn. Thầy không chê bừa bộn trước. Thầy chỉ nói: ``Bàn học gọn không làm em thành thiên tài. Nhưng nó giúp đầu óc em bớt rối. Gọn bên ngoài là một cách tập cho tư duy rõ bên trong.''}
\textit{(The teacher passed a desk covered in scraps and messy books. He didn't criticize the mess directly. He only said: ``A tidy desk won't make you a genius. But it helps your mind become less cluttered. Order outside is one way to train clarity inside.'')}
\end{truyen}

\begin{baihoc}
Sự ngăn nắp tuy nhỏ nhưng giúp việc học nhẹ và rõ hơn nhiều.
\textit{(Orderliness may seem small, but it makes learning much lighter and clearer.)}
\end{baihoc}

\begin{ghinhoanh}
\textbf{Short English to remember:}

A clear desk helps a clear mind.

\end{ghinhoanh}

\begin{tuhocnhanh}
1. Không gian học của em đang giúp hay cản đầu óc em? \textit{Is your study space helping or blocking your mind?}

2. Em có thể dọn gọn điều gì ngay hôm nay? \textit{What can you tidy up today?}
\end{tuhocnhanh}
\ngancach

\section{Muốn tự do lâu dài, phải học kỷ luật từ sớm.}
\begin{truyen}{Tự do và kỷ luật}{Freedom Needs Discipline}
\chuhoa{N}{hiều em thích nghe chữ ``tự do''. Thầy nói: ``Ai cũng muốn tự do. Nhưng tự do lâu dài không đứng trên buông thả. Nó đứng trên kỷ luật. Em biết tự quản mình thì sau này mới có nhiều quyền chọn hơn.''}
\textit{(Many students like the word ``freedom''. The teacher said: ``Everyone wants freedom. But lasting freedom is not built on carelessness. It is built on discipline. If you can manage yourself, later you will have more choices.'')}
\end{truyen}

\begin{baihoc}
Kỷ luật hôm nay mở rộng tự do của ngày mai.
\textit{(Today's discipline widens tomorrow's freedom.)}
\end{baihoc}

\begin{ghinhoanh}
\textbf{Short English to remember:}

Discipline today creates freedom tomorrow.

\end{ghinhoanh}

\begin{tuhocnhanh}
1. Vì sao buông thả không phải là tự do thật? \textit{Why is carelessness not true freedom?}

2. Kỷ luật nào sẽ cho em nhiều lựa chọn hơn sau này? \textit{Which discipline will give you more choices later?}
\end{tuhocnhanh}
\ngancach